\section{Introduction}
\label{sec:introduction}

Huanglongbing (HLB), or citrus greening, is widely regarded as the most destructive disease affecting citrus worldwide, caused by the phloem-restricted bacterium \textit{Candidatus} Liberibacter spp.\ and transmitted primarily by the Asian citrus psyllid \textit{Diaphorina citri} \citep{Bove2006,Halbert2004}. Because no curative treatment exists for the pathogen once a tree is infected, management relies almost entirely on suppressing the vector, and the ectoparasitoid \textit{Tamarixia radiata} is the natural enemy most widely deployed in classical and augmentative biological control programmes against \textit{D.\ citri} \citep{Hall2008, chen2014}. We recently introduced, and analysed at the scale of a single orchard, a four-compartment model that couples the psyllid, the flush and vigour dynamics of its host plant, and \textit{T.\ radiata} under periodic augmentative releases \citep{RodriguezFalcon2026}. That single-patch model showed that the flush--vigour feedback generates a sustained, endogenous oscillation that periodic releases interact with rather than entrain, that sustained releases outperform a single introduction, and it derived an invasion-exponent criterion separating release schedules that merely suppress the pest from those predicted to drive it toward local eradication.

That analysis, however, treated the agroecosystem as a single, well-mixed patch, an idealization shared by most of the theoretical literature on augmentative releases, from the early spatial simulation models developed for boll-weevil parasitoids \citep{ARCASIM1996} to recent optimal-release studies for egg-parasitoid programs against stem borers \citep{Diatraea2024}. Citrus, however, is grown -- and HLB spreads -- across orchards, groves and regions connected by psyllid flight, wind-assisted dispersal and the movement of nursery material, so that a release programme designed for one patch is in practice deployed on a network of patches that continuously exchange both the pest and its natural enemy. Spatial structure is not a secondary detail in host--parasitoid systems, since the spatial arrangement of refuges and the connectivity of the landscape shape host--parasitoid persistence and the suppression that can be achieved \citep{RefugeConnectivity2022}. How precisely a release is placed relative to the host patch measurably changes realized parasitism \citep{PatchDispersal2026}. Despite this, explicit network treatments of augmentative release strategies remain scarce. As recently noted for the sterile insect technique, metapopulation models are commonplace in epidemiology but their application to pest control, and in particular to how a release strategy should be chosen for a given dispersal network, is still limited \citep{SITMetapopulation2026}.

The relevant theoretical toolkit, however, already exists in a closely related field. Representing a spatially extended population as patches (nodes) coupled by dispersal along the edges of a graph is the standard metapopulation formulation used to study epidemic spread \citep{PastorSatorras2001,Colizza2007}, and that literature has repeatedly shown that network topology -- and, in particular, the degree of the node where a perturbation originates -- governs whether it remains confined or invades the whole network \citep{Colizza2007}. Highly connected nodes can act as reservoirs that sustain and re-seed the rest of the network, a role long documented for structured epidemic and metapopulation systems \citep{Colizza2007} and recently shown to matter even for purely ecological, non-epidemic dynamics. Embedding a Lotka--Volterra
community in a network of diffusively coupled patches, rather than treating it as an isolated or merely migration-open system, qualitatively changes which macroecological patterns the model is
able to reproduce \citep{GutierrezArroyo2026}. This body of work suggests that the move from a single patch to a network is not a cosmetic generalization of our earlier model, but can change which outcomes are reachable at all and raises, for an augmentative pest--parasitoid system, three questions that have clear analogues in network epidemiology but had not, to our knowledge, been examined together for this kind of release-driven biological control: does the degree of the outbreak patch control the outcome, as the reservoir-node intuition from epidemic spreading would suggest; does the resulting spatial coherence follow the linear stability of the network's synchronous state, as captured by the master stability function \citep{pecora1998}, or does it instead reflect the strongly nonlinear, finite-time dynamics of an actual outbreak; and how do dispersal and -- realistically -- mortality incurred while moving between patches interact with the choice between releasing uniformly across the network and targeting the outbreak patch alone.

Here we extend the single-patch model of \citet{RodriguezFalcon2026} to an explicit network of $n$ patches, coupling psyllids and parasitoids by diffusion along the edges of seven graph families that span regular, random and scale-free topologies, while flush and plant vigour remain local to each patch, as in the original formulation. One patch is seeded with an outbreak, and augmentative releases of \textit{T.\ radiata} are applied either uniformly across the network or targeted at the outbreak patch alone. The latter acts as the network analogue of the population-wide versus targeted interventions long contrasted in epidemic control on networks \citep{PastorSatorras2001,Colizza2007}. We quantify the outcome with the network-wide pest load and with two efficiency measures that compare strategies using different numbers of released parasitoids, and we complement the simulations with a master-stability analysis of the uniform-release synchronous state. Generalizing the model to a network also required revisiting, and in places correcting, the analysis scripts inherited from the single-patch study. None of the conclusions of that study changes, and we report the corrections together with the new network-level results.

The remainder of the paper is organized as follows. Section \ref{sec:methods} extends the single-node model to a network with diffusive dispersal and transit mortality, defines the response and efficiency variables used to compare release strategies and describes the simulation experiments, which span seven network topologies and a wide range of dispersal and mortality rates, together with the master-stability analysis of the synchronous state. Section \ref{sec:results} reports the results, which include the relative roles of release strategy, topology and the degree of the outbreak patch; the agreement, and disagreement, between the predicted and the observed spatial coherence; and the emergence, once transit mortality is included, of an optimal parasitoid dispersal rate for targeted releases that is absent without it. Section \ref{sec:discussion} discusses the implications for the design of spatially explicit augmentative-release programmes and situates the network results relative to the single-patch model of \citet{RodriguezFalcon2026}.

